\chapter*{\lblHowToUse}
\phantomsection
\addcontentsline{toc}{chapter}{\lblHowToUse}
\markboth{\lblHowToUse}{\lblHowToUse}

\noindent
This book is laid out for a screen and not for a shelf. The page is A4, the
document is one-sided, and every design decision below assumes the same
picture: \textbf{the PDF open on one half of the screen, an editor open on the
other}, and a terminal somewhere under the editor. If you read it on a train
with no laptop you will get the argument and miss the point, because the point
is what happens in the editor.

\section*{The loop}

Every chapter is read the same way, and the loop is short enough to become a
habit by chapter three:

\begin{enumerate}[leftmargin=1.6em, itemsep=3pt]
  \item Read a section. It ends on a listing, and the listing is a
        \emph{file}: its path is printed above it, relative to the repository
        root, and it is the file you open in the editor.
  \item Run the file. Every listing in this book runs as printed, on the
        versions on the title page, and the repository's build runs it again on
        every change. If it does not run for you, one of you has the wrong
        version, and the title page says which.
  \item Do the exercise. An exercise is a starter file and a test beside it.
        Make the test pass, then read what the test checked and why.
  \item Move on. The next section assumes you did the previous three.
\end{enumerate}

\noindent
\mermaidfig{reading-loop}{The loop: read a section in the PDF, open and run its listing in the editor, make the exercise's test pass, move on}{the reading loop}

Here is the whole loop, once, before there is anything to learn from it. The
listing below is a file; open it.

\pyfile{code/ch00/loop.py}{The first listing, and the shape every later one has}{lst:ch00-loop}

\noindent
Run it from the repository root and it prints this, which is also generated
rather than typed \dash{} the line below is what the pinned interpreter printed
when the book was last built:

\transcript{ch00-loop}

\noindent
And the exercise that goes with it. It is numbered 0.1 because this chapter
is not a chapter; the first real one is Exercise~1.1.

\exercise{e00_01_hello}{Make the first test pass}{%
Open the starter file. It contains one function that raises
\code{NotImplementedError}, and the test beside it fails until the function
returns what the test expects. Make it pass. Then run the test again with
\code{-v} and read what pytest printed, because the way pytest reports a
failure is the first thing in this book that has no .NET equivalent worth
naming.}

\noindent
Every exercise in the book prints those three lines: the file you open, the
file that decides whether you are done, and the one command that runs it. The
starter for every exercise fails its test on purpose, and the repository's
build checks that it does \dash{} an exercise whose starter already passes is
not an exercise.

\section*{What a chapter is}

Fourteen chapters, and each is held to the same contract:

\begin{itemize}[leftmargin=1.4em, itemsep=2pt]
  \item \textbf{One argument}, stated in the first paragraph and paid off by
        the last. A chapter is a claim about where a C\# habit stops working
        and what replaces it, not a tour of a library.
  \item \textbf{Under three thousand words} of prose, listings not counted.
        The build counts them and refuses a chapter that is over. Three
        thousand words is a sitting; the book is fourteen sittings.
  \item \textbf{A mapping table}, C\# on the left and Python on the right,
        because the reader already has the concept and needs the name.
  \item \textbf{Listings that are files}, every one of them run on the pinned
        versions before it is printed. Where a listing has not been run it
        says so, in a box, and the build counts the boxes.
  \item \textbf{Traps}, elicited before they are named. Where a C\# habit
        produces Python that runs and is wrong, the chapter lets you write it
        first. Appendix~\ref{app:traps} is the catalogue.
  \item \textbf{Exercises} with a starter and a test, as above.
  \item \textbf{Independence.} A chapter names what it assumes from earlier
        ones and can be read on its own; each was written to be published on
        its own.
\end{itemize}

\section*{The boxes}

\begin{center}
\small
\begin{tabularx}{\linewidth}{@{}lX@{}}
\toprule
\textbf{\lblNote} & Something worth knowing that the argument does not need. \\
\textbf{\lblWarning} & Something that will cost you an afternoon. \\
\textbf{\lblCsharp} & The translation. The C\# side of the sentence you have
  just read, and nothing else; a reader without .NET behind them can skip
  every one of these and lose no content. \\
\textbf{\lblTrap} & A misconception, named after you have walked into it.
  Every one is in Appendix~\ref{app:traps}. \\
\textbf{\lblVersion} & Preview, experimental, or likely to move. Rare in a
  book about a language; common in its last two chapters. \\
\textbf{\lblVerify} & A listing that has not been executed against the pinned
  versions. Nothing ships to a reader with one of these attached. \\
\textbf{\lblExercise} & A starter file and a test. \\
\textbf{\lblProject} & The stage of \code{trace-assert} this section builds.
  See the Introduction. \\
\bottomrule
\end{tabularx}
\end{center}

\section*{Two rules this book holds itself to}

\begin{note}
\textbf{Every listing was run, or says it was not.} The code under
\code{code/} is a project of its own, with a lockfile, and the repository's
build installs exactly the versions on the title page and runs every listing
and every exercise test on every change. A listing that has not been through
that carries a \emph{\lblVerify} box, and the build prints how many there
are. When that number is zero the book is entitled to its title page.
\end{note}

\begin{note}
\textbf{What has not been measured is labelled as judgement.} Where the book
says one thing is faster than another, there is a script under
\code{code/measure/} that produced the number, and the number reaches the page
mechanically rather than by being typed. Where no script has run yet, the
sentence says so, and the table that would hold the numbers stays empty.
Empty tables are not an oversight; they are the mechanism.
\end{note}

\section*{Three releases}

The book ships in three parts, and each is complete on its own.
Version~0.1 is Parts~I and~II \dash{} the runtime, the toolchain and the
language. Version~0.2 adds Parts~III and~IV \dash{} concurrency and
shipping. Version~1.0 adds Part~V and the appendices. A chapter that is not
yet written prints a red box saying so, and the build counts those too.

\section*{If you already write Python}

Read Appendix~\ref{app:traps} first, and be honest about the entries you did
not know. Then read Chapters~\ref{ch:asyncio} and~\ref{ch:testing} anyway,
because they are where a C\# background produces Python that runs and is
wrong, and a year of writing Python is not always enough to notice.

\section*{Both languages, one source}

This book exists in English and in Polish, built from one repository. The two
editions share their chapters, their listings, their exercises and their
numbering; only the prose differs. A check on every build proves the two have
the same structure and the same numbers, so a correction applied to one
edition cannot silently miss the other. English is the first language of the
book: the code, the identifiers and every listing comment are English in both
editions, because the reader is being trained to read English code.
Where a Polish term has no settled form in real use, the Polish edition
keeps the English word and inflects it, and says so where it does.
